\documentclass[11pt]{article}
\usepackage[a4paper,margin=26mm]{geometry}
\usepackage{amsmath,amssymb}
\usepackage[hidelinks]{hyperref}
\title{Conjecture 00000001042: an exact affine variance and the fixed-degree asymptotic obstruction}
\author{lizaixi01, with OpenAI Codex assistance}
\date{22 October 2026}
\begin{document}
\maketitle

\section{Statement, interpretation, and scope}
The source conjecture says: ``The variance of the value-set size of a random
polynomial $f$ uniform of degree $\leq d$ is asymptotically
$q(1-1/e)/d$.'' The Chinese version likewise specifies degree at most $d$.
The source is conjecture 00000001042 at upstream commit:\par\noindent
\texttt{b355e42906e6eaf52bc7e9d005e0ef02d8aa7400}.

We disprove the reading in which $q$ tends to infinity through finite-field
cardinalities while $d$ is fixed, with the stated class including $d=1$.
The source does not specify a limiting regime. We do not claim to settle a
different formulation restricted to growing $d$, or one excluding $d=1$.
Uniform means uniform on all coefficient vectors, including zero leading
coefficients and the zero polynomial. In degree at most one this is the
uniform distribution on the $q^2$ distinct formal polynomials $aX+b$.

\section{Exact computation for every finite field}
Let $K$ be any finite field of cardinality $q$, and let $(a,b)$ be uniform
on $K^2$. Define the actual value-set size
\[
 V(a,b)=\#\{ax+b:x\in K\}.
\]
If $a=0$ the image is the singleton $\{b\}$, so $V=1$.
If $a\ne0$, the inverse image of $y$ is $(y-b)/a$; hence the map is
surjective and $V=q$. There are $q$ pairs with zero slope and $q(q-1)$
with nonzero slope. Thus
\[
 \mathbb E[V^k]=\frac{1+(q-1)q^k}{q},\qquad
 \mathbb E[V]=q-1+\frac1q.
\]
The variance, equivalently the uniform mean squared deviation, is
\[
 \operatorname{Var}(V)=\mathbb E[V^2]-\mathbb E[V]^2
     =\frac{(q-1)^3}{q^2}.                         \tag{1}
\]
This is an identity for every finite field, not a finite numerical test.

\section{Asymptotic contradiction}
Let $p_n$ enumerate the primes in increasing order, and take
$K_n=\mathbb Z/p_n\mathbb Z$. These are finite fields and $p_n\to\infty$.
Equation (1) gives
\[
 \frac{\operatorname{Var}_{K_n}(V)}{p_n}
    =\left(1-\frac1{p_n}\right)^3\longrightarrow 1.
\]
But $\exp(1)>1$, so $0<1-1/\exp(1)<1$.
The proposed fixed-$d=1$ equivalence would instead imply a normalized
limit $1-1/\exp(1)$. Uniqueness of limits yields a contradiction.
In Lean we additionally refute the exact ratio formulation
\[
 \frac{\operatorname{Var}_{K_n}(V)}{p_n(1-1/\exp(1))}\longrightarrow 1.
\]
Failure along this legitimate unbounded sequence already disproves the
asserted $q$-asymptotic reading.

\section{Formalization and semantic correspondence}
The standalone project uses Lean 4.33.0 and the pinned Mathlib revision in the project configuration and dependency manifest.
The main theorems in \texttt{lean/Main.lean} are:
\begin{itemize}
\item \texttt{affineEquiv}: an explicit bijection between coefficient pairs
and all formal polynomials of natural degree at most one;
\texttt{affine\_eval} connects their evaluations to $ax+b$.
Thus the uniform coefficient enumeration is exactly the requested uniform
polynomial ensemble.
\item \texttt{valueCount\_zero}, \texttt{valueCount\_ne\_zero}, and
\texttt{sum\_power}: actual finite image cardinalities and exact sums.
\item \texttt{variance\_is\_mean\_squared\_deviation} and
\texttt{variance\_formula}: the standard variance definition and (1), with
no assumed moment or counting formula.
\item \texttt{normalized\_variance\_limit}: the true limit along actual
prime fields. The final theorem
\texttt{refutes\_asymptotic\_equivalence} proves failure of the
conjectured ratio limit.
\end{itemize}
The only foundational axioms allowed are Lean's usual
\texttt{propext}, \texttt{Classical.choice}, and \texttt{Quot.sound}.
There is no admitted proof, native decision shortcut, or new axiom.
Reproduction commands and exact verification receipts accompany this report.

\section{Previous submission and what is repaired}
\href{https://github.com/The-Last-Math-Competition/The-Last-Math-Competition/pull/155}{PR 155}
already presented this elementary mathematical obstruction. It was closed
unmerged: its Lean proof checked only $q=101$, while the general identity
and the limiting argument remained in prose. A finite discrepancy alone
does not disprove an asymptotic claim. The
\href{https://github.com/The-Last-Math-Competition/The-Last-Math-Competition/pull/155\#issuecomment-5951812972}{maintainer's review}
identified both missing formal steps and the unspecified limiting regime.
This submission credits that prior observation, proves the general
finite-field identity and the limiting contradiction in Lean, and explicitly
limits its conclusion to the fixed-degree, growing-field interpretation.
It makes no claim to new mathematics or to resolving every possible
alternative interpretation of the incomplete source statement.

\section{Provenance}
OpenAI Codex performed the reasoning, Lean formalization, and report
preparation for the submitting account. A separate supervising Codex agent
read the original statement and full proof source for semantic review;
this is an AI review, not a human certification. The official competition
review and acceptance remain the maintainers' responsibility.
\end{document}
